% Preambula terpadu untuk tujuh akar mandiri catatan teori representasi Duncan.
% Berkas akar terjemahan tidak diubah; docmute mengabaikan preambula lokalnya.
\usepackage[indonesian]{babel}
\usepackage{fontspec}
\defaultfontfeatures{Ligatures=TeX,Scale=MatchLowercase}
\setmainfont{TeX Gyre Termes}
\setsansfont{TeX Gyre Heros}
\setmonofont{Latin Modern Mono}

\usepackage{amssymb,amsmath,amsthm}
\usepackage{ytableau}
\usepackage{geometry}
\geometry{a4paper,margin=28mm,headheight=16pt,footskip=30pt}
\usepackage{microtype}
\usepackage{csquotes}
\DeclareQuoteAlias{english}{indonesian}
\usepackage{docmute}
\usepackage{fancyhdr}
\usepackage{lastpage}
\usepackage[
  backend=biber,
  style=alphabetic,
  sorting=nyt,
  giveninits=true,
  maxbibnames=99,
  maxcitenames=3
]{biblatex}
\DeclareLanguageMapping{indonesian}{english}
\addbibresource{../source/rep_theory.bib}
\usepackage{xurl}
\usepackage[
  unicode,
  colorlinks=false,
  pdfborder={0 0 0},
  pdftitle={Catatan Teori Representasi - adaptasi bahasa Indonesia},
  pdfauthor={Alexander Duncan},
  pdfsubject={Adaptasi bahasa Indonesia dari catatan MATH 742},
  pdfkeywords={teori representasi, aljabar, grup hingga, CC BY 4.0},
  pdfcreator={OpenAI Codex gpt-5.6-sol, Ultra}
]{hyperref}

% Gaya teorema terpadu. Semua lingkungan bernomor memakai pencacah yang sama.
\theoremstyle{plain}
\newtheorem{theorem}{Teorema}[section]
\newtheorem{lemma}[theorem]{Lema}
\newtheorem{proposition}[theorem]{Proposisi}
\newtheorem{corollary}[theorem]{Korolari}
\newtheorem{conjecture}[theorem]{Konjektur}
\newtheorem{exercise}[theorem]{Latihan}
\theoremstyle{definition}
\newtheorem{definition}[theorem]{Definisi}
\newtheorem{example}[theorem]{Contoh}
\theoremstyle{remark}
\newtheorem{remark}[theorem]{Catatan}
\numberwithin{equation}{section}

% Identitas edisi dan provenance yang tidak berubah antarbangunan.
\newcommand{\DuncanSourceAuthor}{Alexander Duncan}
\newcommand{\DuncanSourceCourse}{MATH 742, Musim Semi 2023}
\newcommand{\DuncanSourceCommit}{c62d36f41189da4bd3da4671668f68720df54ff7}
\newcommand{\DuncanSourceTree}{e83ee440666133b14dec440158a108069a13e9e4}
\newcommand{\DuncanSourceRepository}{https://github.com/vtorsor/representation-theory-notes}
\newcommand{\DuncanLicenseName}{Creative Commons Attribution 4.0 International}
\newcommand{\DuncanLicenseSPDX}{CC-BY-4.0}
\newcommand{\DuncanLicenseURL}{https://creativecommons.org/licenses/by/4.0/}
\newcommand{\DuncanModelProvenance}{OpenAI Codex gpt-5.6-sol, Ultra, atas instruksi pengguna}
\newcommand{\DuncanEditionDate}{Edisi sumber: Musim Semi 2023}

\pagestyle{fancy}
\fancyhf{}
\lhead{MATH 742}
\chead{Catatan Teori Representasi}
\rhead{Adaptasi id-ID}
\lfoot{\DuncanEditionDate}
\rfoot{\thepage\ dari \pageref{LastPage}}
\renewcommand{\footrulewidth}{0.4pt}
\setlength{\parindent}{1.5em}
\setlength{\parskip}{0pt}
\setlength{\emergencystretch}{3em}
\date{\DuncanEditionDate}

% \maketitle pada kelas article menonaktifkan dirinya secara global.
% Simpan definisi awal agar setiap akar mandiri tetap memperoleh halaman judul.
\let\DuncanArticleMakeTitle\maketitle
